%
% einleitung.tex -- Beispiel-File für die Einleitung
%
% (c) 2020 Prof Dr Andreas Müller, Hochschule Rapperswil
%
% !TEX root = ../../buch.tex
% !TEX encoding = UTF-8
%
\section{Fakultät erweitern\label{gamma:section:teil0}}
\kopfrechts{Fakultät erweitern}

Die Fakultät ist definiert als
\begin{equation}
n! = \prod_{k=1}^{n} k = 1 \times 2 \times 3 \times \cdots \times n
\label{gamma:fakultaet}
\end{equation}
und gilt für die natürlichen Zahlen. Sie fällt vorallem durch ihr schnelles Wachstum auf.
Nicht ganz intuitiv dabei ist, dass $0! = 1$ ist. Stellt man die Fakultät in die rekursive Form um
\begin{equation}
n! = n \times (n-1)! <=> (n-1) = n!/n
\label{gamma:fakultaet-rekursiv}
\end{equation}
und setzte 0 für $(n-1)$ ein, zeigt sich warum $0! = 1!/1 = 1$ ist. 

Eine erste Definition der Fakultät für nicht Natürliche Zahlen wird bereits 1720 von Daniel Bernoulli und Christian Goldbach gegeben. 
Auch Leonhard Euler und James Stirling steuern im Laufe der Jahrhunderte ihre eigene Definition bei. 
Man sieht, die Erweiterung der Fakultät hat grosse Namen der Mathematik beschäftigt. 

\subsection{Eigenschaften der Fakultät}

Um die Fakultät zu erweitern, müssen drei essenzielle Eigenschaften festgelegt werden. 
\begin{enumerate}
  \item $f(1) = 1$
  \item $f(x+1) = x \times f(x), x > 0$
  \item $f(x)$ wächst schnell
\end{enumerate}
Erfüllt eine Funktion alle diese Eigenschaften, kann sie als Erweiterung der Fakultät für nicht Natürliche Zahlen betrachtet werden.


